\section{Penutup}

Bab ini berisi ringkasan hasil yang diperoleh serta saran untuk pengembangan
lebih lanjut.

\subsection{Kesimpulan}

\begin{enumerate}
	\item \textbf{Judol Detector berhasil menerapkan algoritma pencocokan string}
	(KMP, Boyer-Moore, Aho-Corasick, Rabin-Karp, RegEx, dan Weighted Levenshtein)
	untuk mendeteksi indikasi konten judol pada halaman web.
	\item \textbf{Algoritma RegEx dan exact matching memberikan performa terbaik}
	untuk halaman umum, sedangkan Weighted Levenshtein efektif menangkap
	variasi visual namun memiliki biaya komputasi tinggi.
	\item \textbf{Hasil pengujian menunjukkan kebutuhan pengendalian false positive,}
	terutama pada pola RegEx yang dapat menangkap token umum tanpa konteks.
\end{enumerate}

\subsection{Saran}

\begin{enumerate}
	\item \textbf{Tambahkan filter konteks untuk RegEx} misalnya pengecualian
	prefix umum atau verifikasi dengan keyword sekitar agar false positive
	berkurang.
	\item \textbf{Optimasi Weighted Levenshtein dengan pembatasan kandidat} dan
	pengaturan threshold adaptif agar performa tetap responsif pada halaman besar.
\end{enumerate}
